\documentclass[10pt]{article}
\usepackage[T1]{fontenc}
\usepackage[utf8]{inputenc}
\usepackage[french]{babel}
\usepackage{amsmath,amssymb,mathrsfs,mathpazo}
\usepackage{geometry}
\geometry{a4paper,margin=2cm}
\usepackage{graphicx,array,tabularx,tabularray,colortbl,multirow,multicol,booktabs,enumitem,pifont,ulem,fancybox,diagbox,cancel}
\usepackage{pstricks,pstricks-add,pst-plot,pst-tree,pst-node,pst-text,pst-func,pst-math,pst-bspline,pst-3dplot,pst-eucl,pst-solides3d}
\let\pscalculate\undefined
\usepackage{pst-bezier}
\usepackage{tikz,pgfplots,tkz-tab}
\usetikzlibrary{arrows,arrows.meta,patterns,calc,positioning,trees,decorations.markings,shapes.geometric,shapes.misc,shapes.arrows,intersections,decorations.pathmorphing,decorations.pathreplacing,backgrounds,shadows,babel}
\pgfplotsset{compat=1.18}
\usepgfplotslibrary{fillbetween}
\usepackage[np]{numprint}
\usepackage{listings,makecell,pas-tableur,twoopt}
\let\barre\undefined
\usepackage[pstricks]{tablvar}
\definecolor{darkgreen}{RGB}{0,100,0}
\providecommand{\red}{\color{red}}
\providecommand{\blue}{\color{blue}}
\providecommand{\green}{\color{green}}
\providecommand{\black}{\color{black}}
\providecommand{\white}{\color{white}}
\providecommand{\gray}{\color{gray}}
\providecommand{\cyan}{\color{cyan}}
\providecommand{\magenta}{\color{magenta}}
\providecommand{\yellow}{\color{yellow}}
\providecommand{\violet}{\color{violet}}
\providecommand{\orange}{\color{orange}}
\providecommand{\euro}{\texteuro}
\providecommand{\EUR}{\texteuro}
\providecommand{\ieme}{\textsuperscript{e}}
\providecommand{\degres}{\ensuremath{^\circ}}
\providecommand{\degree}{\ensuremath{^\circ}}
\providecommand{\vv}[1]{\overrightarrow{#1}}
\providecommand{\Coord}[1]{#1}
\providecommand{\decosix}{}
\providecommand{\starredbullet}{\ensuremath{\bullet}}
\pagestyle{empty}
\setlength{\parindent}{0pt}
\usepackage{adjustbox}\DeclareUnicodeCharacter{2236}{:}\DeclareUnicodeCharacter{2190}{\ensuremath{\leftarrow}}
\providecommand{\R}{\mathbb{R}}
\providecommand{\N}{\mathbb{N}}
\providecommand{\Z}{\mathbb{Z}}
\providecommand{\Q}{\mathbb{Q}}
\providecommand{\D}{\mathbb{D}}
\providecommand{\C}{\mathbb{C}}
\providecommand{\e}{{\mathrm{e}}}
\providecommand{\vect}[1]{\overrightarrow{#1}}
\providecommand{\vectt}[1]{\overrightarrow{\mathrm{#1}}}
\providecommand{\vvt}[1]{\overrightarrow{\mathrm{#1}}}
\providecommand{\np}[2][]{{#2~#1}}
\providecommand{\vv}[1]{{\overrightarrow{#1}}}
\renewcommand{\d}{{\mathrm{d}}}
\providecommand{\degres}{{^\circ}}
\providecommand{\degre}{{^\circ}}
\providecommand{\blue}{\color{blue}}
\providecommand{\red}{\color{red}}
\providecommand{\enskip}{\quad}
\providecommand{\llbracket}{\lbrack\!\lbrack}
\providecommand{\rrbracket}{\rbrack\!\rbrack}
\providecommand{\euro}{{\text{€}}}
\providecommand{\no}{{\text{n°}}}
\providecommand{\interval}[2]{{[#1;#2]}}
\providecommand{\intervFO}[2]{{[#1;#2[}}
\providecommand{\dotfill}{\ldots}

\providecommand{\ds}{\displaystyle}
\providecommand{\textding}[1]{#1}
\providecommand{\Oij}{\(\mathrm{O};\vec{\imath},\vec{\jmath}\)}
\providecommand{\Oijk}{\(\mathrm{O};\vec{\imath},\vec{\jmath},\vec{k}\)}
\providecommand{\Ouv}{\(\mathrm{O};\vec{u},\vec{v}\)}
\providecommand{\intervFF}[2]{\ensuremath{[#1;#2]}}
\renewcommand{\euro}{\ifmmode\text{\texteuro}\else\texteuro\fi}
\DeclareUnicodeCharacter{00B0}{\ifmmode^{\circ}\else\textdegree\fi}
\DeclareUnicodeCharacter{20AC}{\ifmmode\text{\texteuro}\else\texteuro\fi}
% Énoncés : Éducation nationale. Transcription et corrigés : APMEP, Denis Vergès et contributeurs. Découpage et classement : CoopMaths, Sébastien Lozano, Éric Elter et contributeurs. Adaptation de la présentation : Maths Explorer.
% CC BY-SA : https://coopmaths.fr/www/about/
% https://coopmaths.fr/alea/static/bac/2026/tex/bac_2026_06_sujet2_polynesie_2.tex
\begin{document}
\noindent\textit{APMEP / CoopMaths. CC BY-SA. Présentation adaptée par Maths Explorer.}\par\medskip
Dans cet exercice, les parties A et B sont indépendantes.

\subsection*{Partie A}

On considère la fonction $f$ définie pour tout $x \in [4~;~10]$ par $f(x) = \ln(8x^3 - 1)$.

On admet que la fonction $f$ est dérivable sur $[4~;~10]$ et on note $f'$ sa fonction dérivée.

\medskip

\begin{enumerate}
\item Montrer pour tout $x \in [4~;~10]$, on a :

\[f'(x) = \frac{24x^2}{(2x-1)(4x^2+2x+1)}.\]

\item[2.] 
	\begin{enumerate}[label=\alph*.]
		\item En déduire que $f$ est croissante sur $[4~;~10]$.
		\item Montrer que, pour tout $x \in [4~;~10]$, on a :

\[4 \leqslant f(x) \leqslant 10.\]

	\end{enumerate}
\item[3.] On définit la fonction $g$ définie pour tout $x \in [4~;~10]$ par $g(x) = f(x) - x$.

On admet que $g$ est strictement décroissante sur son ensemble de définition.

\medskip

\begin{enumerate}[label=\alph*.]
\item Montrer que, sur l'intervalle $[4~;~10]$, l'équation $g(x) = 0$ admet une unique solution $\alpha$.
\item Déterminer un encadrement à $10^{-3}$ près de $\alpha$.
    \end{enumerate}
\end{enumerate}

\subsection*{Partie B}

On admet qu'on peut définir la suite $(u_n)$ par $u_0 = 5$ et pour tout entier naturel $n$, par $u_{n+1} = f(u_n)$, où $f$ est la fonction étudiée à la partie A.

\medskip

\begin{enumerate}
\item Donner une valeur approchée de $u_1$ à $10^{-1}$ près.
\item[2.] 
	\begin{enumerate}[label=\alph*.]
		\item Montrer, à l'aide d'un raisonnement par récurrence que, pour tout $n \in \mathbb{N}$, on a :
		
\[4 \leqslant u_n \leqslant u_{n+1} \leqslant 10.\]
		\item En déduire que la suite $(u_n)$ converge.
	\end{enumerate}
\item[3.] Montrer que la suite $(u_n)$ converge vers le nombre $\alpha$, défini dans la partie A.
\end{enumerate}

\subsection*{Partie C}

Dans cette partie, on admet que $8{,}4 \leqslant \alpha \leqslant 8{,}5$.

On a créé, en langage Python, la fonction \texttt{alpha} suivante.

On précise qu'en langage Python, l'instruction \texttt{log} correspond à $\ln$.


\begin{center}
\begin{tabular}{|l|}\hline
\texttt{def f(x) :}\\
\quad \texttt{return log(8*x**3 - 1)}\\
~\\
\texttt{def alpha(a):}\\
\quad \texttt{n = 0}\\
\quad \texttt{u = 5}\\
\quad \texttt{while u < a:}\\
\quad \quad \texttt{u = f(u)}\\
\quad \quad \texttt{n = n + 1}\\
\quad \texttt{return n}\\ \hline
\end{tabular}
\end{center}


\medskip

\begin{enumerate}
\item On admet que \texttt{alpha(8.499)} renvoie la valeur 11.

Interpréter ce résultat dans le contexte de l'exercice.
\item Un utilisateur exécute \texttt{alpha(8.6)}.

Expliquer pourquoi aucune valeur n'est alors renvoyée.
\end{enumerate}

\bigskip

\end{document}
